\hypertarget{chapter-03-section-04-page-231}{}
\section{用半离散给出的另一存在性证明}
\label{sec:ch3-semidiscrete-existence}
本节给出 Navier--Stokes 方程弱解的另一种存在性证明, 它适用于任意空间维数. 通过关于 $t$ 的半离散构造逼近解, 再利用紧性论证取极限. 第 \ref{subsec:ch3-semidiscrete-formulation} 小节以适用于任意维数的方式重新表述问题并陈述存在性结果; 第 \ref{subsec:ch3-semidiscrete-approximation} 小节构造逼近解; 第 \ref{subsec:ch3-semidiscrete-estimates} 和 \ref{subsec:ch3-semidiscrete-limit} 小节分别处理先验估计与极限过程.

\subsection{问题的陈述}
\label{subsec:ch3-semidiscrete-formulation}
在更高维数中陈述存在性定理之前, 必须重新表述弱解问题. 与定常情形一样, 当 $n>4$ 时, $b$ 在 $V$ 上不是连续三线性型, 因而像 \eqref{eq:ch3-existence-force-assumption}--\eqref{eq:ch3-existence-weak-initial-condition} 那样的陈述没有意义, 因为其中的 $b(u(t),u(t),v)$ 项可能没有定义.

为此再次引入(见第 \ref{chap:steady-state-navier-stokes} 章第 \ref{subsec:ch2-sobolev-compactness} 小节)空间 $V_s$:
\begin{equation}
V_s=\overline{\mathcal V}^{\,\mathbf H_0^1(\Omega)\cap\mathbf H^s(\Omega)},
\qquad s\geq1.
\label{eq:ch3-semidiscrete-vs-definition}
\end{equation}

\hypertarget{chapter-03-page-232}{}
空间 $\mathbf H_0^1(\Omega)\cap\mathbf H^s(\Omega)$ 和 $V_s$ 赋予 $\mathbf H^s(\Omega)$ 的通常 Hilbert 范数:
\begin{equation}
\lVert u\rVert_{\mathbf H^s(\Omega)}=
\left\{\sum_{|j|\leq s}|D^ju|^2\right\}^{1/2}
\qquad(s\text{ 为整数}).
\label{eq:ch3-semidiscrete-hs-norm}
\end{equation}
显然, 对 $s\geq1$,
\begin{equation}
V_s\subset V
\label{eq:ch3-semidiscrete-vs-embedding}
\end{equation}
为连续注入, 且 $V_s$ 在 $V$ 中稠密. 当 $s\geq n/2$ 时, $b$ 定义在 $V\times V\times V_s$ 上; 更确切地说:

\MBSourceStatementNumber{lemma}{1}
\begin{Lemma}
\label{lem:ch3-semidiscrete-trilinear-continuity}
若 $s\geq n/2$, 则 $b$ 在 $V\times V\times V_s$ 上连续三线性, 且
\begin{equation}
|b(u,v,w)|\leq c|u|\lVert v\rVert\lVert w\rVert_{V_s}.
\label{eq:ch3-semidiscrete-trilinear-estimate}
\end{equation}
\end{Lemma}

\begin{Proof}
对 $u,v,w\in\mathcal V$, Hölder 不等式给出
\[
\begin{aligned}
|b(u,v,w)|=|b(u,w,v)|
&\leq\sum_{i,j=1}^n\lVert u_i\rVert_{L^2(\Omega)}
\lVert D_iw_j\rVert_{L^n(\Omega)}
\lVert v_j\rVert_{L^{2n/(n-2)}(\Omega)}\\
&\leq c_0|u|\lVert v\rVert\sum_{i,j=1}^n
\lVert D_iw_j\rVert_{L^n(\Omega)},
\end{aligned}
\]
这里使用了 Sobolev 不等式 $H_0^1(\Omega)\subset L^{2n/(n-2)}(\Omega)$. 因 $s\geq n/2$, $H^{s-1}(\Omega)$ 包含于 $L^q(\Omega)$, 其中
\begin{equation}
\frac1q=\frac12-\frac{s-1}{n},
\qquad q\geq n.
\label{eq:ch3-semidiscrete-sobolev-exponent}
\end{equation}
若 $w\in V_s$, 则 $D_iw_j\in H^{s-1}(\Omega)\cap L^2(\Omega)$, 因而 $D_iw_j\in L^n(\Omega)$ 且
\[
\lVert D_iw_j\rVert_{L^n(\Omega)}\leq c_1\lVert w\rVert_{V_s}.
\]
所以
\begin{equation}
|b(u,v,w)|\leq c_2|u|\lVert v\rVert\lVert w\rVert_{V_s}.
\label{eq:ch3-semidiscrete-trilinear-estimate-dense}
\end{equation}
此估计把 $b$ 从 $\mathcal V\times\mathcal V\times\mathcal V$ 连续延拓到 $V\times V\times V_s$, 甚至延拓到 $H\times V\times V_s$.
\end{Proof}

\MBSourceStatementNumber{lemma}{2}
\begin{Lemma}
\label{lem:ch3-semidiscrete-bu-vs-dual}
若 $u\in L^2(0,T;V)\cap L^\infty(0,T;H)$, 则对 $s\geq n/2$, $Bu\in L^2(0,T;V_s')$.
\end{Lemma}

\begin{Proof}
由 $B$ 的定义和 \eqref{eq:ch3-semidiscrete-trilinear-estimate},
\[
|\langle Bu(t),v\rangle|=|b(u(t),u(t),v)|
\leq c|u(t)|\lVert u(t)\rVert\lVert v\rVert_{V_s},
\qquad\forall v\in V_s.
\]
故对几乎每个 $t\in[0,T]$,
\[
\lVert Bu(t)\rVert_{V_s'}\leq c|u(t)|\lVert u(t)\rVert,
\]
结论成立.
\end{Proof}

\hypertarget{chapter-03-page-233}{}
在任意空间维数中, Navier--Stokes 问题可作如下弱表述.

\MBProblemHeading
\label{prob:ch3-semidiscrete-variational}
给定 $f,u_0$ 满足
\begin{equation}
f\in L^2(0,T;V'),
\label{eq:ch3-semidiscrete-force}
\end{equation}
\begin{equation}
u_0\in H,
\label{eq:ch3-semidiscrete-initial-data}
\end{equation}
求 $u$ 满足
\begin{equation}
u\in L^2(0,T;V)\cap L^\infty(0,T;H),
\label{eq:ch3-semidiscrete-solution-space}
\end{equation}
\begin{equation}
\frac d{dt}(u,v)+\nu((u,v))+b(u,u,v)=\langle f,v\rangle,
\qquad\forall v\in V_s,quad s\geq\frac n2,
\label{eq:ch3-semidiscrete-variational-equation}
\end{equation}
\begin{equation}
u(0)=u_0.
\label{eq:ch3-semidiscrete-initial-condition}
\end{equation}
若 $u$ 满足 \eqref{eq:ch3-semidiscrete-solution-space} 和 \eqref{eq:ch3-semidiscrete-variational-equation}, 则
\[
\frac d{dt}\langle u,v\rangle=\langle g,v\rangle,
\qquad\forall v\in V_s,
\]
其中 $g=f-Bu-\nu Au$. 由引理 \ref{lem:ch3-semidiscrete-bu-vs-dual}, $Bu\in L^2(0,T;V_s')$, 而 $f-\nu Au\in L^2(0,T;V')$, 故
\begin{equation}
g\in L^2(0,T;V_s').
\label{eq:ch3-semidiscrete-g-space}
\end{equation}
引理 \ref{lem:ch3-linear-banach-valued-derivative} 因而推出
\begin{equation}
u'\in L^2(0,T;V_s'),
\qquad u'=f-\nu Au-Bu.
\label{eq:ch3-semidiscrete-time-derivative}
\end{equation}
于是 $u$ 几乎处处等于从 $[0,T]$ 到 $V_s'$ 的连续函数, \eqref{eq:ch3-semidiscrete-initial-condition} 有意义.

问题 \ref{prob:ch3-semidiscrete-variational} 的另一表述如下.

\MBProblemHeading
\label{prob:ch3-semidiscrete-operator}
给定满足 \eqref{eq:ch3-semidiscrete-force}--\eqref{eq:ch3-semidiscrete-initial-data} 的 $f,u_0$, 求 $u$ 满足
\begin{equation}
u\in L^2(0,T;V)\cap L^\infty(0,T;H),
\qquad u'\in L^2(0,T;V_s'),quad s\geq\frac n2,
\label{eq:ch3-semidiscrete-operator-space}
\end{equation}
\begin{equation}
u'+\nu Au+Bu=f\quad\text{于 }(0,T),
\label{eq:ch3-semidiscrete-operator-equation}
\end{equation}
\begin{equation}
u(0)=u_0.
\label{eq:ch3-semidiscrete-operator-initial}
\end{equation}
两种表述等价. 下列定理给出解的存在性, 并蕴含定理 \ref{thm:ch3-existence-weak-solution}.

\MBSourceStatementNumber{theorem}{1}
\begin{Theorem}
\label{thm:ch3-semidiscrete-existence}
给定满足 \eqref{eq:ch3-semidiscrete-force}--\eqref{eq:ch3-semidiscrete-initial-data} 的 $f,u_0$. 则问题 \ref{prob:ch3-semidiscrete-operator} 至少存在一个解 $u$. 此外, $u$ 从 $[0,T]$ 到 $H$ 弱连续.
\end{Theorem}

该定理在第 \ref{subsec:ch3-semidiscrete-approximation} 和 \ref{subsec:ch3-semidiscrete-estimates} 小节证明; $H$ 中的弱连续性直接来自 \eqref{eq:ch3-semidiscrete-operator-space} 和引理 \ref{lem:ch3-linear-weak-continuity-transfer}.

\hypertarget{chapter-03-page-234}{}
\subsection{逼近解}
\label{subsec:ch3-semidiscrete-approximation}
令 $N$ 为稍后趋于无穷的整数, 并置
\begin{equation}
k=\frac TN.
\label{eq:ch3-semidiscrete-time-step}
\end{equation}
递归定义 $V$ 中元素 $u^0,u^1,\ldots,u^N$; $u^m$ 在某种意义下逼近区间 $mk<t<(m+1)k$ 上所求函数 $u$.

先定义 $V'$ 中元素 $f^1,\ldots,f^N$:
\begin{equation}
f^m=\frac1k\int_{(m-1)k}^{mk}f(t)\,dt,
\qquad m=1,\ldots,N;quad f^m\in V'.
\label{eq:ch3-semidiscrete-force-average}
\end{equation}
从
\begin{equation}
u^0=u_0
\label{eq:ch3-semidiscrete-u0}
\end{equation}
开始; 已知 $u^0,\ldots,u^{m-1}$ 时, 把 $u^m$ 定义为满足
\begin{equation}
\frac{u^m-u^{m-1}}k+\nu Au^m+Bu^m=f^m
\label{eq:ch3-semidiscrete-recursion}
\end{equation}
的 $V$ 中元素. $u^m$ 依赖于 $k$, 但简记为 $u^m$ 而不写 $u_k^m$. 其存在性由引理 \ref{lem:ch3-semidiscrete-step-existence} 保证.

\MBSourceStatementNumber{lemma}{3}
\begin{Lemma}
\label{lem:ch3-semidiscrete-step-existence}
对每个固定的 $k$ 和每个 $m\geq1$, 至少存在一个满足 \eqref{eq:ch3-semidiscrete-recursion} 的 $u^m$, 且
\begin{equation}
|u^m|^2-|u^{m-1}|^2+|u^m-u^{m-1}|^2
+2k\nu\lVert u^m\rVert^2\leq2k\langle f^m,u^m\rangle.
\label{eq:ch3-semidiscrete-step-energy}
\end{equation}
\end{Lemma}

对每个固定的 $k$(或 $N$), 由 $u^1,\ldots,u^N$ 定义逼近函数:
\begin{equation}
u_k:[0,T]\to V,
\qquad u_k(t)=u^m,quad t\in[(m-1)k,mk],\quad m=1,\ldots,N,
\label{eq:ch3-semidiscrete-piecewise-constant}
\end{equation}
\begin{equation}
w_k:[0,T]\to H,
\quad w_k\text{ 连续且在每个 }[(m-1)k,mk]\text{ 上线性},
\quad w_k(mk)=u^m,quad m=0,\ldots,N.
\label{eq:ch3-semidiscrete-piecewise-linear}
\end{equation}
第 \ref{subsec:ch3-semidiscrete-estimates} 小节给出这些函数的先验估计; 然后令 $k\to0$ 取极限(源文写作 \MBUnresolvedReference{section}{Section 4.3}).

\begin{Proof}[引理 \ref{lem:ch3-semidiscrete-step-existence} 的证明]
等式 \eqref{eq:ch3-semidiscrete-recursion} 必须在比 $V'$ 更大的空间中理解, 例如 $V_s'$ 中, $s\geq n/2$. 它等价于
\begin{equation}
(u^m,v)+k\nu((u^m,v))+kb(u^m,u^m,v)
=\langle u^{m-1}+kf^m,v\rangle,
\qquad\forall v\in V_s.
\label{eq:ch3-semidiscrete-step-variational}
\end{equation}
基本像定理 \MBUnresolvedReference{theorem}{Theorem 2.1.2} 一样使用 Galerkin 方法. 选取 $V_s$ 中自由且完备, 因而在 $V$ 中也完备的序列 $w_1,\ldots,w_i,\ldots$. 对每个 $r$, 应用引理 \MBUnresolvedReference{lemma}{Lemma 1.4} 得到
\begin{equation}
\varphi_r=\sum_{i=1}^r\xi_{i,r}w_i,
\label{eq:ch3-semidiscrete-galerkin-phi}
\end{equation}
满足
\begin{equation}
(\varphi_r,v)+k\nu((\varphi_r,v))+kb(\varphi_r,\varphi_r,v)
=\langle u^{m-1}+kf^m,v\rangle,
\quad\forall v\in\operatorname{Sp}[w_1,\ldots,w_r].
\label{eq:ch3-semidiscrete-galerkin-equation}
\end{equation}

\hypertarget{chapter-03-page-235}{}
还需取得与 $r$ 无关的先验估计并令 $r\to\infty$(此证明中 $k,m$ 固定). 在 \eqref{eq:ch3-semidiscrete-galerkin-equation} 中取 $v=\varphi_r$ 得
\begin{equation}
(\varphi_r-u^{m-1},\varphi_r)+k\nu\lVert\varphi_r\rVert^2
=k\langle f^m,\varphi_r\rangle.
\label{eq:ch3-semidiscrete-galerkin-energy}
\end{equation}
又有
\begin{equation}
2(a-b,a)=|a|^2-|b|^2+|a-b|^2,
\qquad\forall a,b\in H,
\label{eq:ch3-semidiscrete-polarization}
\end{equation}
故
\begin{equation}
\begin{aligned}
|\varphi_r|^2+|\varphi_r-u^{m-1}|^2+2k\nu\lVert\varphi_r\rVert^2
&=|u^{m-1}|^2+2k\langle f^m,\varphi_r\rangle\\
&\leq|u^{m-1}|^2+k\nu\lVert\varphi_r\rVert^2
+\frac k\nu\lVert f^m\rVert_{V'}^2.
\end{aligned}
\label{eq:ch3-semidiscrete-galerkin-energy-bound}
\end{equation}
因此
\begin{equation}
|\varphi_r|^2+|\varphi_r-u^{m-1}|^2+k\nu\lVert\varphi_r\rVert^2
\leq|u^{m-1}|^2+\frac k\nu\lVert f^m\rVert_{V'}^2.
\label{eq:ch3-semidiscrete-galerkin-uniform-bound}
\end{equation}
于是 $\varphi_r$ 在 $V$ 中有界, 可抽取子序列 $\varphi_{r'}$ 使
\begin{equation}
\varphi_{r'}\rightharpoonup\varphi\quad\text{在 }V\text{ 中},
\qquad r'\to\infty.
\label{eq:ch3-semidiscrete-galerkin-weak-limit}
\end{equation}
用标准论证在 \eqref{eq:ch3-semidiscrete-galerkin-equation} 中取极限, 证明 $\varphi=u^m$ 满足 \eqref{eq:ch3-semidiscrete-step-variational}. 最后在 \eqref{eq:ch3-semidiscrete-galerkin-energy-bound} 中取下极限, 利用弱拓扑下范数的下半连续性, 得到 \eqref{eq:ch3-semidiscrete-step-energy}.
\end{Proof}

\subsection{先验估计}
\label{subsec:ch3-semidiscrete-estimates}
\MBSourceStatementNumber{lemma}{4}
\begin{Lemma}
\label{lem:ch3-semidiscrete-discrete-estimates}
有
\begin{equation}
|u^m|^2\leq d_1,
\qquad m=1,\ldots,N,
\label{eq:ch3-semidiscrete-um-bound}
\end{equation}
\begin{equation}
k\sum_{m=1}^N\lVert u^m\rVert^2\leq\frac{d_1}\nu,
\label{eq:ch3-semidiscrete-v-sum-bound}
\end{equation}
\begin{equation}
\sum_{m=1}^N|u^m-u^{m-1}|^2\leq d_1,
\label{eq:ch3-semidiscrete-increment-bound}
\end{equation}
其中 $d_1$ 只依赖于数据:
\begin{equation}
d_1=|u_0|^2+\frac1\nu\int_0^T\lVert f(s)\rVert_{V'}^2\,ds.
\label{eq:ch3-semidiscrete-d1}
\end{equation}
\end{Lemma}

\begin{Proof}
如引理 \ref{lem:ch3-semidiscrete-step-existence} 的证明所述, 至少当 $n>4$ 时不能在 \eqref{eq:ch3-semidiscrete-recursion} 中与 $u^m$ 作内积, 因为 $Bu^m$ 一般不属于 $V'$. 但 \eqref{eq:ch3-semidiscrete-step-energy} 起到同样作用.

\hypertarget{chapter-03-page-236}{}
以
\[
2k\lVert f^m\rVert_{V'}\lVert u^m\rVert
\leq k\nu\lVert u^m\rVert^2+\frac k\nu\lVert f^m\rVert_{V'}^2
\]
控制右端, 得
\begin{equation}
|u^m|^2-|u^{m-1}|^2+|u^m-u^{m-1}|^2+k\nu\lVert u^m\rVert^2
\leq\frac k\nu\lVert f^m\rVert_{V'}^2.
\label{eq:ch3-semidiscrete-step-energy-reduced}
\end{equation}
对 $m=1,\ldots,N$ 求和得
\begin{equation}
|u^N|^2+\sum_{m=1}^N|u^m-u^{m-1}|^2
+k\nu\sum_{m=1}^N\lVert u^m\rVert^2
\leq|u_0|^2+\frac k\nu\sum_{m=1}^N\lVert f^m\rVert_{V'}^2.
\label{eq:ch3-semidiscrete-global-energy}
\end{equation}
对 $m=1,\ldots,r$ 求和并舍去 $\lVert u^m\rVert^2$ 项, 得
\begin{equation}
|u^r|^2\leq|u_0|^2+\frac k\nu\sum_{m=1}^r\lVert f^m\rVert_{V'}^2
\leq|u_0|^2+\frac k\nu\sum_{m=1}^N\lVert f^m\rVert_{V'}^2,
\quad r=1,\ldots,N.
\label{eq:ch3-semidiscrete-partial-energy}
\end{equation}
结合下一引理即可得证.
\end{Proof}

\MBSourceStatementNumber{lemma}{5}
\begin{Lemma}
\label{lem:ch3-semidiscrete-force-average-bound}
若 $f^m$ 由 \eqref{eq:ch3-semidiscrete-force-average} 定义, 则
\begin{equation}
k\sum_{m=1}^N\lVert f^m\rVert_{V'}^2
\leq\int_0^T\lVert f(t)\rVert_{V'}^2\,dt.
\label{eq:ch3-semidiscrete-force-average-estimate}
\end{equation}
\end{Lemma}

\begin{Proof}
由 Schwarz 不等式,
\[
\lVert f^m\rVert_{V'}^2
=\left\lVert\frac1k\int_{(m-1)k}^{mk}f(t)\,dt\right\rVert_{V'}^2
\leq\frac1k\int_{(m-1)k}^{mk}\lVert f(t)\rVert_{V'}^2\,dt.
\]
对 $m=1,\ldots,N$ 求和即得 \eqref{eq:ch3-semidiscrete-force-average-estimate}.
\end{Proof}

\MBSourceStatementNumber{lemma}{6}
\begin{Lemma}
\label{lem:ch3-semidiscrete-time-difference-bound}
和式
\[
k\sum_{m=1}^N
\left\lVert\frac{u^m-u^{m-1}}k\right\rVert_{V_s'}^2
\]
与 $k$ 无关地有界.
\end{Lemma}

\begin{Proof}
在 \eqref{eq:ch3-semidiscrete-recursion} 中取 $V_s'$ 范数, 得
\[
\begin{aligned}
\left\lVert\frac{u^m-u^{m-1}}k\right\rVert_{V_s'}
&\leq\lVert f^m\rVert_{V_s'}+\nu\lVert Au^m\rVert_{V_s'}
+\lVert Bu^m\rVert_{V_s'}\\
&\leq c_1(\lVert f^m\rVert_{V'}+\lVert u^m\rVert)+\lVert Bu^m\rVert_{V_s'},
\end{aligned}
\]
所以
\[
\left\lVert\frac{u^m-u^{m-1}}k\right\rVert_{V_s'}^2
\leq c_2\{\lVert f^m\rVert_{V'}^2+\lVert u^m\rVert^2
+\lVert Bu^m\rVert_{V_s'}^2\}.
\]
由 \eqref{eq:ch3-semidiscrete-trilinear-estimate} 和 \eqref{eq:ch3-semidiscrete-um-bound},
\[
\lVert Bu^m\rVert_{V_s'}^2\leq c_3|u^m|^2\lVert u^m\rVert^2
\leq c_4\lVert u^m\rVert^2.
\]

\hypertarget{chapter-03-page-237}{}
最终
\[
k\sum_{m=1}^N
\left\lVert\frac{u^m-u^{m-1}}k\right\rVert_{V_s'}^2
\leq c_5k\sum_{m=1}^N(\lVert f^m\rVert_{V'}^2+\lVert u^m\rVert^2),
\]
结合 \eqref{eq:ch3-semidiscrete-v-sum-bound} 和引理 \ref{lem:ch3-semidiscrete-force-average-bound} 即得证.
\end{Proof}

现在把上述估计解释为逼近函数的估计.

\MBSourceStatementNumber{lemma}{7}
\begin{Lemma}
\label{lem:ch3-semidiscrete-function-bounds}
$u_k,w_k$ 在 $L^2(0,T;V)\cap L^\infty(0,T;H)$ 的一个有界集中; $w_k'$ 在 $L^2(0,T;V_s')$ 中有界, 且
\begin{equation}
u_k-w_k\longrightarrow0\quad\text{在 }L^2(0,T;H)\text{ 中},
\qquad k\to0.
\label{eq:ch3-semidiscrete-uk-wk-convergence}
\end{equation}
\end{Lemma}

\begin{Proof}
关于 $u_k,w_k,w_k'$ 的估计分别是 \eqref{eq:ch3-semidiscrete-um-bound}--\eqref{eq:ch3-semidiscrete-v-sum-bound} 与引理 \ref{lem:ch3-semidiscrete-time-difference-bound} 的解释; \eqref{eq:ch3-semidiscrete-uk-wk-convergence} 由 \eqref{eq:ch3-semidiscrete-increment-bound} 和下一引理推出.
\end{Proof}

\MBSourceStatementNumber{lemma}{8}
\begin{Lemma}
\label{lem:ch3-semidiscrete-interpolation-distance}
有
\begin{equation}
|u_k-w_k|_{L^2(0,T;H)}
=\left\{\frac k3\sum_{m=1}^N|u^m-u^{m-1}|^2\right\}^{1/2}.
\label{eq:ch3-semidiscrete-interpolation-distance-formula}
\end{equation}
\end{Lemma}

\begin{Proof}
当 $(m-1)k\leq t\leq mk$ 时,
\[
w_k(t)-u_k(t)=\frac{t-mk}{k}(u^m-u^{m-1}),
\]
且
\[
\int_{(m-1)k}^{mk}|w_k(t)-u_k(t)|^2\,dt
=\frac k3|u^m-u^{m-1}|^2.
\]
求和即得 \eqref{eq:ch3-semidiscrete-interpolation-distance-formula}.
\end{Proof}

\subsection{取极限}
\label{subsec:ch3-semidiscrete-limit}
由引理 \ref{lem:ch3-semidiscrete-function-bounds}, 可从 $u_k$ 抽取子序列 $u_{k'}$, 使
\begin{equation}
u_{k'}\rightharpoonup u\quad\text{在 }L^2(0,T;V)\text{ 中},
\qquad
u_{k'}\overset{*}{\rightharpoonup}u\quad\text{在 }L^\infty(0,T;H)\text{ 中}.
\label{eq:ch3-semidiscrete-uk-weak-limits}
\end{equation}
为在 \eqref{eq:ch3-semidiscrete-recursion} 中取极限, 需要 $u_k$ 的强收敛; $w_k$ 将发挥辅助作用. 可选取子序列使
\begin{equation}
w_{k'}\rightharpoonup u_*\quad\text{在 }L^2(0,T;V)\text{ 中},
\qquad
w_{k'}\overset{*}{\rightharpoonup}u_*quad\text{在 }L^\infty(0,T;H)\text{ 中},
\label{eq:ch3-semidiscrete-wk-weak-limits}
\end{equation}
\begin{equation}
\frac{dw_{k'}}{dt}\rightharpoonup u_*'
\quad\text{在 }L^2(0,T;V_s')\text{ 中}.
\label{eq:ch3-semidiscrete-wk-derivative-limit}
\end{equation}
由 \eqref{eq:ch3-semidiscrete-uk-wk-convergence}, $u_*=u$. 定理 \ref{thm:ch3-compactness-aubin-lions-banach} 给出
\begin{equation}
w_{k'}\longrightarrow u\quad\text{在 }L^2(0,T;H)\text{ 中},
\label{eq:ch3-semidiscrete-wk-strong-limit}
\end{equation}
再由 \eqref{eq:ch3-semidiscrete-uk-wk-convergence},
\begin{equation}
u_{k'}\longrightarrow u\quad\text{在 }L^2(0,T;H)\text{ 中}.
\label{eq:ch3-semidiscrete-uk-strong-limit}
\end{equation}

\hypertarget{chapter-03-page-238}{}
方程 \eqref{eq:ch3-semidiscrete-recursion} 可解释为
\begin{equation}
\frac{dw_k}{dt}+Au_k+Bu_k=f_k,
\label{eq:ch3-semidiscrete-functional-equation}
\end{equation}
其中
\[
f_k(t)=f^m,
\qquad(m-1)k\leq t\leq mk,quad m=1,\ldots,N.
\]
由 \eqref{eq:ch3-semidiscrete-uk-weak-limits}, \eqref{eq:ch3-semidiscrete-uk-strong-limit} 及引理 \ref{lem:ch3-existence-bu-integrability}, \ref{lem:ch3-semidiscrete-bu-vs-dual},
\[
Bu_{k'}\rightharpoonup Bu\quad\text{在 }L^2(0,T;V_s')\text{ 中}.
\]
由下面的引理 \ref{lem:ch3-semidiscrete-force-convergence},
\[
f_k\longrightarrow f\quad\text{在 }L^2(0,T;V')\text{ 中}.
\]
因此可在 \eqref{eq:ch3-semidiscrete-functional-equation} 中取极限, 得
\[
u'+\nu Au+Bu=f.
\]
由 \eqref{eq:ch3-semidiscrete-wk-weak-limits}, \eqref{eq:ch3-semidiscrete-wk-derivative-limit} 和引理 \ref{lem:ch3-semidiscrete-trilinear-continuity},
\[
\langle w_{k'}(t),\sigma\rangle\longrightarrow\langle u(t),\sigma\rangle,
\qquad\forall\sigma\in V_s',\quad\forall t\in[0,T].
\]
因 $w_{k'}(0)=u_0$, 得 $u(0)=u_0$. 这证明 $u$ 满足 \eqref{eq:ch3-semidiscrete-operator-space}--\eqref{eq:ch3-semidiscrete-operator-initial}; 证明下列引理后定理 \ref{thm:ch3-semidiscrete-existence} 即证完.

\MBSourceStatementNumber{lemma}{9}
\begin{Lemma}
\label{lem:ch3-semidiscrete-force-convergence}
\begin{equation}
f_k\longrightarrow f\quad\text{在 }L^2(0,T;V')\text{ 中},
\qquad k\to0.
\label{eq:ch3-semidiscrete-force-convergence}
\end{equation}
\end{Lemma}

\begin{Proof}
变换 $f\mapsto f_k$ 是 $L^2(0,T;V')$ 中的线性平均映射; 由引理 \ref{lem:ch3-semidiscrete-force-average-bound}, 它连续且
\begin{equation}
\lVert f_k\rVert_{L^2(0,T;V')}
\leq\lVert f\rVert_{L^2(0,T;V')}.
\label{eq:ch3-semidiscrete-averaging-contraction}
\end{equation}
因此只需在 $L^2(0,T;V')$ 的稠密子空间中证明 \eqref{eq:ch3-semidiscrete-force-convergence}; 对 $f\in C([0,T];V')$, 结论显然, 略去证明.
\end{Proof}

\MBSourceStatementNumber{remark}{1}
\begin{Remark}
\label{rem:ch3-semidiscrete-energy-inequality}
将 \eqref{eq:ch3-semidiscrete-step-energy} 对 $m=1,\ldots,r$ 求和并舍去 $|u^m-u^{m-1}|^2$ 项, 得
\begin{equation}
|u^r|^2+2k\nu\sum_{m=1}^r\lVert u^m\rVert^2
\leq|u_0|^2+2k\sum_{m=1}^r\langle f^m,u^m\rangle.
\label{eq:ch3-semidiscrete-discrete-energy-remark}
\end{equation}
它可解释为
\begin{equation}
|u_k(t)|^2+2\nu\int_0^{t_k}\lVert u_k(s)\rVert^2\,ds
\leq|u_0|^2+\int_0^{t_k}\langle f_k(s),u_k(s)\rangle\,ds,
\label{eq:ch3-semidiscrete-functional-energy-remark}
\end{equation}
其中
\begin{equation}
t(k)=(m+k)k,
\qquad mk\leq t<(m+1)k.
\label{eq:ch3-semidiscrete-time-ceiling-source}
\end{equation}
\end{Remark}
